%! Unit 1: Linear Functions — Practice Test (KEY)
% GENERATED BODY: the blank and key bodies are byte-identical except the header title;
% answers live in \slot / \optok / work, which the preamble shows or hides.
\documentclass[12pt]{article}
\usepackage{algebra2-article}
\usepackage{algebra2-key}   % pulls in -boxes; do NOT also load -boxes
% Coordinate grid: {xmin}{xmax}{ymin}{ymax}
\newcommand{\lgrid}[4]{%
  \draw[step=1, linegray!40, very thin] (#1,#3) grid (#2,#4);
  \draw[->, semithick, charcoal] (#1,0)--(#2,0) node[below right, font=\scriptsize]{$x$};
  \draw[->, semithick, charcoal] (0,#3)--(0,#4) node[left, font=\scriptsize]{$y$};}
% Endpoint dots: {x}{y}. Closed = filled; open = white-filled ring.
\newcommand{\fdot}[2]{\fill[forest] (#1,#2) circle (4.5pt);}
\newcommand{\hdot}[2]{\draw[very thick, forest, fill=white] (#1,#2) circle (4.5pt);}
% Part-divider strip (headlinebox takes one arg: the background color)
\newcommand{\parthead}[1]{%
  \boxguard[10]\vspace{6pt}\noindent
  \begin{headlinebox}{forest}\color{white}\bfseries #1\end{headlinebox}%
  \vspace{4pt}}
% THE WORK RULE for a test: \slot{width}{answer} and \opt / \optok are the ONLY
% lines that differ between blank and key, and they differ in the preamble, not the body.
% A slot is a fixed-width underline in both files, so no line rewraps in the key.
\newcommand{\slot}[2]{\underline{\makebox[#1]{\ans{#2}}}}
\newcommand{\opt}[2]{(#1)~#2}
\newcommand{\optok}[2]{\mbox{\llap{\textcolor{keyred}{$\boldsymbol{\rightarrow}$}\,}\textcolor{keyred}{\textbf{(#1)}~#2}}}

\begin{document}

\pageheader{Unit 1: Linear Functions}{Practice Test --- Answer Key}
\namedateperiod

\vspace{0.06in}
\noindent\textbf{Instructions:} Show all work. Write every line in $y=mx+b$ form when asked.
No notes or calculator unless your teacher says so. \hfill\textbf{Total: 100 pts}

\vspace{4pt}
\boxguard[6]
\begin{remindbox}
\small
\textbf{This is a practice test.} It mirrors the real Unit 1 test in length, format, and the
ideas it covers, but the numbers differ. Work every problem with no notes, then check your answers
against the key.
\end{remindbox}

% ======================================================================
\parthead{Part A --- Types of Slope (8 pts)}
For each line, write whether its slope is \textbf{positive}, \textbf{negative}, \textbf{zero}, or
\textbf{undefined}.

\vspace{6pt}
\noindent
\begin{minipage}[t]{0.24\textwidth}\centering
\begin{tikzpicture}[scale=0.34]
  \lgrid{-3}{3}{-3}{3}
  \draw[very thick, navy] (-3,2)--(3,-1);
\end{tikzpicture}\\[2pt]
1.\ \slot{2.4cm}{negative}
\end{minipage}%
\begin{minipage}[t]{0.24\textwidth}\centering
\begin{tikzpicture}[scale=0.34]
  \lgrid{-3}{3}{-3}{3}
  \draw[very thick, navy] (2,-3)--(2,3);
\end{tikzpicture}\\[2pt]
2.\ \slot{2.4cm}{undefined}
\end{minipage}%
\begin{minipage}[t]{0.24\textwidth}\centering
\begin{tikzpicture}[scale=0.34]
  \lgrid{-3}{3}{-3}{3}
  \draw[very thick, navy] (-1,-3)--(2,3);
\end{tikzpicture}\\[2pt]
3.\ \slot{2.4cm}{positive}
\end{minipage}%
\begin{minipage}[t]{0.24\textwidth}\centering
\begin{tikzpicture}[scale=0.34]
  \lgrid{-3}{3}{-3}{3}
  \draw[very thick, navy] (-3,-1.5)--(3,-1.5);
\end{tikzpicture}\\[2pt]
4.\ \slot{2.4cm}{zero}
\end{minipage}

\vspace{10pt}
\noindent No graph this time --- decide from the points or the equation.
\begin{enumerate}[leftmargin=1.9em, itemsep=6pt, label=\arabic*., start=5]
  \item The line through $(3,\,2)$ and $(-1,\,2)$: \hfill \slot{2.4cm}{zero}
  \item The line $y=5x-1$: \hfill \slot{2.4cm}{positive}
  \item The line $x=-4$: \hfill \slot{2.4cm}{undefined}
  \item The line through $(0,\,6)$ and $(2,\,1)$: \hfill \slot{2.4cm}{negative}
\end{enumerate}

% ======================================================================
\parthead{Part B --- Writing Equations of Lines (20 pts)}
Write each line in slope-intercept form, $y=mx+b$. Show how you found $m$ and $b$.

\begin{enumerate}[leftmargin=1.9em, itemsep=10pt, label=\arabic*., start=9]
  \item The line through $(1,\,5)$ and $(3,\,9)$.
    \begin{work}
      m &= \frac{9-5}{3-1} = \frac{4}{2} = 2 \\
      5 &= 2(1) + b \quad\Rightarrow\quad b = 3
    \end{work}
    \vspace{0.5cm}
    \hfill Equation: \slot{3.6cm}{$y=2x+3$}

  \item The line through $(-4,\,5)$ and $(2,\,-4)$.
    \begin{work}
      m &= \frac{-4-5}{2-(-4)} = \frac{-9}{6} = -\frac{3}{2} \\
      5 &= -\frac{3}{2}(-4) + b = 6 + b \quad\Rightarrow\quad b = -1
    \end{work}
    \vspace{0.5cm}
    \hfill Equation: \slot{3.6cm}{$y=-\tfrac{3}{2}x-1$}

  \item The line with slope $2$ that passes through $(3,\,4)$.
    \begin{work}
      4 &= 2(3) + b \\
      4 &= 6 + b \quad\Rightarrow\quad b = -2
    \end{work}
    \vspace{0.5cm}
    \hfill Equation: \slot{3.6cm}{$y=2x-2$}

  \item The line with slope $-\tfrac{1}{3}$ that passes through $(6,\,1)$.
    \begin{work}
      1 &= -\frac{1}{3}(6) + b \\
      1 &= -2 + b \quad\Rightarrow\quad b = 3
    \end{work}
    \vspace{0.5cm}
    \hfill Equation: \slot{3.6cm}{$y=-\tfrac{1}{3}x+3$}
\end{enumerate}

% ======================================================================
\parthead{Part C --- Piecewise Functions (20 pts)}

\begin{enumerate}[leftmargin=1.9em, itemsep=10pt, label=\arabic*., start=13]
  \item \textbf{Evaluate.} (8 pts) \quad
    $f(x)=\begin{cases} 3x+2, & x<-1 \\ -1, & -1\le x\le 1 \\ 2x-5, & x>1 \end{cases}$
    \\[8pt]
    $f(-2)=$ \slot{1.4cm}{$-4$} \qquad $f(-1)=$ \slot{1.4cm}{$-1$} \qquad
    $f(1)=$ \slot{1.4cm}{$-1$} \qquad $f(4)=$ \slot{1.4cm}{$3$}

  \item \textbf{Domain.} (12 pts) Give the domain of each function as an inequality or in interval
    notation.
    \begin{enumerate}[label=(\alph*), leftmargin=1.9em, itemsep=12pt]
      \item $g(x)=\begin{cases} x-2, & x<3 \\ 7, & 3\le x<6 \end{cases}$
        \hfill Domain: \slot{3.8cm}{$x<6$, \ $(-\infty,\,6)$}
      \item $h(x)=\begin{cases} x+1, & -2\le x\le 0 \\ 4-x, & 0<x\le 5 \end{cases}$
        \hfill Domain: \slot{3.8cm}{$-2\le x\le 5$, \ $[-2,\,5]$}
      \boxguard[9]
      \item The function $k$ graphed below.
        \\[4pt]
        \begin{minipage}[c]{0.45\linewidth}\centering
        \begin{tikzpicture}[scale=0.5]
          \lgrid{-4}{5}{-3}{4}
          \draw[very thick, forest] (-3,3)--(0,0);
          \fdot{-3}{3} \hdot{0}{0}
          \draw[very thick, forest] (0,2)--(4,2);
          \fdot{0}{2} \hdot{4}{2}
        \end{tikzpicture}
        \end{minipage}%
        \hfill Domain: \slot{3.8cm}{$-3\le x<4$, \ $[-3,\,4)$}
    \end{enumerate}
\end{enumerate}

% ======================================================================
\parthead{Part D --- The Absolute Value Function (28 pts)}

\begin{enumerate}[leftmargin=1.9em, itemsep=10pt, label=\arabic*., start=15]
  \item (3 pts) Circle the letter of the graph of an \textbf{absolute value} function.
    \\[4pt]
    \begin{minipage}[t]{0.24\linewidth}\centering
    \begin{tikzpicture}[scale=0.32]
      \lgrid{-3}{3}{-3}{3}
      \draw[very thick, navy] (-3,2)--(-1,0)--(2,3);
    \end{tikzpicture}\\ \optok{A}{}
    \end{minipage}%
    \begin{minipage}[t]{0.24\linewidth}\centering
    \begin{tikzpicture}[scale=0.32]
      \lgrid{-3}{3}{-3}{3}
      \draw[very thick, navy] (-3,-2.5)--(3,1.5);
    \end{tikzpicture}\\ \opt{B}{}
    \end{minipage}%
    \begin{minipage}[t]{0.24\linewidth}\centering
    \begin{tikzpicture}[scale=0.32]
      \lgrid{-3}{3}{-3}{3}
      \draw[very thick, navy, domain=-2.2:2.2, smooth] plot (\x, {-0.5*\x*\x+1.5});
    \end{tikzpicture}\\ \opt{C}{}
    \end{minipage}%
    \begin{minipage}[t]{0.24\linewidth}\centering
    \begin{tikzpicture}[scale=0.32]
      \lgrid{-3}{3}{-3}{3}
      \foreach \k in {-2,...,1} { \draw[very thick, navy] (\k,\k)--(\k+1,\k);
        \fill[navy] (\k,\k) circle (4pt); \draw[thick, navy, fill=white] (\k+1,\k) circle (4pt); }
    \end{tikzpicture}\\ \opt{D}{}
    \end{minipage}

  \item (3 pts) Which equation is an \textbf{absolute value} function?
    \\[4pt]
    \opt{A}{$y=x^2-3$} \hfill \opt{B}{$y=2x+5$} \hfill \optok{C}{$y=|x+1|-3$} \hfill
    \opt{D}{$y=\lfloor x\rfloor+1$}

  \item (15 pts) Each function is a transformation of the parent $y=|x|$. Fill in the table.
    \\[6pt]
    {\renewcommand{\arraystretch}{1.9}
    \begin{tabular}{|c|c|c|c|c|}
      \hline
      \textbf{Function} & \textbf{Vertex} & \textbf{Left/right shift} & \textbf{Up/down shift} &
      \textbf{Opens} \\ \hline
      $g(x)=|x+2|+3$ & \slot{1.9cm}{$(-2,\,3)$} & \slot{2.1cm}{left 2} & \slot{2.1cm}{up 3} &
        \slot{1.4cm}{up} \\ \hline
      $h(x)=|x-5|-4$ & \slot{1.9cm}{$(5,\,-4)$} & \slot{2.1cm}{right 5} & \slot{2.1cm}{down 4} &
        \slot{1.4cm}{up} \\ \hline
      $k(x)=-|x|-2$ & \slot{1.9cm}{$(0,\,-2)$} & \slot{2.1cm}{none} & \slot{2.1cm}{down 2} &
        \slot{1.4cm}{down} \\ \hline
    \end{tabular}}

  \item (4 pts) Which equation matches the graph?
    \\[4pt]
    \begin{minipage}[c]{0.36\linewidth}\centering
    \begin{tikzpicture}[scale=0.36]
      \lgrid{-1}{7}{-2}{4}
      \draw[very thick, navy] (-0.5,2.5)--(3,-1)--(6.5,2.5);
      \fill[navy] (3,-1) circle (4pt);
    \end{tikzpicture}
    \end{minipage}%
    \begin{minipage}[c]{0.6\linewidth}
    \opt{A}{$y=|x+3|-1$} \\[6pt]
    \opt{B}{$y=|x-3|+1$} \\[6pt]
    \optok{C}{$y=|x-3|-1$} \\[6pt]
    \opt{D}{$y=-|x-3|-1$}
    \end{minipage}

  \item (3 pts) Compared with the parent $y=|x|$, the graph of $y=\tfrac{1}{2}|x|$ is
    \\[4pt]
    \opt{A}{narrower} \hfill \optok{B}{wider} \hfill \opt{C}{shifted down 2} \hfill
    \opt{D}{shifted left 2}
\end{enumerate}

% ======================================================================
\parthead{Part E --- Absolute Value Equations \& Inequalities (24 pts)}
Solve. Show both cases.

\begin{enumerate}[leftmargin=1.9em, itemsep=10pt, label=\arabic*., start=20]
  \item (8 pts) \quad $|2x+1|=9$
    \begin{work}
      2x+1 &= 9 \quad\text{or}\quad 2x+1 = -9 \\
      2x &= 8 \quad\text{or}\quad 2x = -10
    \end{work}
    \vspace{0.5cm}
    \hfill Solution: \slot{4.2cm}{$x=4$ or $x=-5$}

  \item (8 pts) \quad $|x-3|\le 2$
    \begin{work}
      -2 &\le x-3 \le 2 \\
      1 &\le x \le 5
    \end{work}
    \vspace{0.5cm}
    \hfill Solution: \slot{4.2cm}{$1\le x\le 5$}

  \item (8 pts) \quad $|x+4|>6$
    \begin{work}
      x+4 &> 6 \quad\text{or}\quad x+4 < -6 \\
      x &> 2 \quad\text{or}\quad x < -10
    \end{work}
    \vspace{0.5cm}
    \hfill Solution: \slot{4.2cm}{$x>2$ or $x<-10$}
\end{enumerate}

\end{document}
